\documentclass[10pt,a4paper]{amsart}
\usepackage[margin=19mm]{geometry}
\usepackage{amsmath,amssymb,mathtools}
\usepackage[scr=rsfs,cal=euler]{mathalfa}
\usepackage{xfrac}
\usepackage[unicode,hidelinks,bookmarks=false]{hyperref}
\newcommand{\ie}{i.e.\ }
\newcommand{\cf}{cf.\ }
\newcommand{\resp}{respectively\ }
\newcommand{\Subsection}{\subsection}
\providecommand{\nobreakdash}{\nobreak\hbox{-}}
\setlength{\emergencystretch}{2em}
\multlinegap0pt
\usepackage{bm}
\usepackage{bbm}
\usepackage{dsfont}
\usepackage{xspace}
\usepackage{cancel}
\usepackage{mathtools}
\usepackage{amsthm}
\makeatletter
\@ifundefined{theorem}{\newtheorem{theorem}{Theorem}}{}
\@ifundefined{lemma}{\newtheorem{lemma}[theorem]{Lemma}}{}
\makeatother
\newcommand{\D}{\mathrm{D}}
\newcommand{\E}{\mathbb{E}}
\newcommand{\HH}{\mathcal{H}}
\newcommand{\R}{\mathbb{R}}
\newcommand{\U}{\mathcal{U}}
\renewcommand{\O}{\mathcal{O}}
\renewcommand{\SS}{\mathbb{S}}
\title[Structural and Solution Analysis for the Ordered Weber Probl]{Structural and Solution Analysis for the Ordered Weber Problem under Spatial Uncertainty}
\author{V\'ictor Blanco, Miguel Mart\'inez-Ant\'on}
\date{Source: arXiv:2511.01481v2 (CC BY 4.0)}
\begin{document}
\maketitle

\noindent\emph{Source and licence.} Structural and Solution Analysis for the Ordered Weber Problem under Spatial Uncertainty. Version arXiv:2511.01481v2 (CC BY 4.0), distributed under the Creative Commons Attribution 4.0 International licence, as recorded in the arXiv licence metadata for this exact version and shown on the abstract page https://arxiv.org/abs/2511.01481v2. Full rights notice: licenses/arxiv-2511.01481-CC-BY-4.0.txt.\par\smallskip

\noindent\textit{Editorial scope.} Counted unit: the Lemma whose source label is \texttt{lemma:single-va} in the article (source0.tex lines 600--602, source label \texttt{lemma:single-va}), delivered together with the article's own complete original proof of it (source0.tex lines 603--659, one \texttt{proof} environment of 57 lines). No mathematics is written, repaired or re-derived: statement and proof are the article's own text line for line. Required context retained uncounted: eq:owp (lines 200--202). Every definition, display and label of those lines is retained unchanged. Omitted from this package and not redistributed: 1--199, 203--599, 660--1780. Nothing inside a retained excerpt is omitted or altered: every retained body is delivered exactly as the article's lines are written, so no omission receipt of any kind accompanies this package. Citations: the retained excerpts carry no citation command at all, so no bibliography entry of the article is transcribed and no reference is redistributed. Reference adaptation: none was needed and none was made; every reference inside the delivered unit resolves inside it, so no unresolved reference or citation is produced. Printed slips kept literal: no printed word, symbol, hypothesis or number of the counted statement or of its proof was altered. Layout and class: the article's own class is \texttt{amsart} with packages including geometry, multicol, adjustbox, graphicx, multirow, amsmath, amssymb, amsfonts, amsthm, mathrsfs, xcolor, booktabs, natbib, hyperref; the wrapper is the lane's pinned \texttt{amsart} editorial wrapper at 10pt on A4, which restates the theorem-like environments the retained text needs and adds the article's own macro definitions that the retained text needs (\texttt{\textbackslash D}, \texttt{\textbackslash E}, \texttt{\textbackslash HH}, \texttt{\textbackslash O}, \texttt{\textbackslash R}, \texttt{\textbackslash SS}, \texttt{\textbackslash U}), with extra wrapper packages (bm, bbm, dsfont, xspace, cancel, mathtools). The delivered numbering is the unit's own. Assets: no retained span contains a figure environment, a graphics-inclusion command, a PDF-page inclusion, a table or a TikZ picture; no image file is copied into this package, the package declares no asset dependency and no image file is redistributed with it. Licence: version arXiv:2511.01481v2 of this article is distributed under the Creative Commons Attribution 4.0 International licence, as recorded in the arXiv licence metadata for this exact version and shown on the abstract page https://arxiv.org/abs/2511.01481v2. A full rights notice is delivered at licenses/arxiv-2511.01481-CC-BY-4.0.txt. Characters: every retained excerpt is pure ASCII, so the delivered engine, XeLaTeX with its default Latin Modern text font, reports no missing character.
\begin{equation}\label{eq:owp}\tag{\rm OWP-SU}
    \min_{y \in \R^d} \O_{\bm{\lambda}}\left(\omega_{1}\mathbb{E}\left[\D\left(y,X_{1}\right)\right],\ldots, \omega_{n}\mathbb{E}\left[\D\left(y,X_{n}\right)\right]\right).
\end{equation}

\begin{lemma}\label{lemma:single-va}
    Let $(\Omega,\mu)$ be a probability space and $y^*\in \R^d$ be a point. If $\mu$ is a spherically symmetric measure centered in $y^*$, then $y^*$ is optimal solution to \ref{eq:owp} for any $X$ drawn according to $\mu$.
\end{lemma}

\begin{proof}
    It will be sufficient to prove that for every $y\neq y^*$, the inequality $\E\left[\|y-X\|\right]\geq \E\left[\|y^*-X\|\right]$ holds for any $X$ drawn according to $\mu$. First, it begins with the one-dimensional case $d=1$. Let $y\in \R$ be a real, then we have
    \begin{equation*}
        \E\left[|y-X|\right]=\int_{-\infty}^y(y-x)d\mu + \int_{y}^{+\infty}(x-y)d\mu.
    \end{equation*}

    We can assume that $y>y*$. Then, we have
    \begin{equation*}
        \E\left[|y-X|\right] - \E\left[|y^*-X|\right] = \int_{-\infty}^{y^*}(y-y^*)d\mu +\int_{y^*}^{y} (y+y^*-2x) d\mu + \int_{y}^{+\infty} (y^*-y) d\mu.
    \end{equation*}
    Since $\mu$ is invariant under rotations centered in $y^*$, then $\int_{-\infty}^{y^*}(y-y^*)d\mu = \int_{y^*}^{+\infty}(y-y^*)d\mu$. Thus, we get
    \begin{equation*}
        \E\left[|y-X|\right] - \E\left[|y^*-X|\right] = 2\int_{y^*}^{y} (y-x) d\mu \geq 0.
    \end{equation*}
    The same procedure can be made for $y<y^*$.

    Now, we follow with the case $d\geq 2$. Note that we can write any $y\in \R^d$ as $y = ru+y^*$, for a scalar $r\geq 0$ and a unit vector $u$. Since $\mu$ is invariant under rotations centered in $y^*$, to prove the claim it suffices to show that for any fixed unit vector $u\in \SS^{d-1}_1(0)$, it follows that $r = 0$ is the minimizer of $\E\left[\|ru+y^*-X\|\right]$ in $\R_+$. Notice that $\E\left[\|ru+y^*-X\|\right]$ is a continuous function in $r$, since for every $\epsilon>0$ and for every nonnegative $r$ and $r'$ such that $|r'-r|<\epsilon$, we have
    \begin{align*}
        \left|\E\left[\|ru+y^*-X\|\right] - \E\left[\|r'u+y^*-X\|\right]\right|&=\left|\E\left[\|ru+y^*-X\|-\|r'u+y^*-X\|\right]\right|\\
        & \leq \|r'u-ru\|=|r'-r|<\epsilon.
    \end{align*}
    Hence, by the Newton-Leibniz formula, we have

    \begin{equation*}
        \E\left[\|su+y^*-X\|\right]-\E\left[\|y^*-X\|\right]=\int_0^s\frac{d}{dr}\E\left[\|ru+y^*-X\|\right]dr, \; \forall s>0.
    \end{equation*}
    Thus, to see the result, it is sufficient to show that
    \begin{equation*}
        \frac{d}{dr}\E\left[\|ru+y^*-X\|\right]>0, \; \forall r>0.
    \end{equation*}
    
    We know that
    \begin{align}
        \frac{d}{dr}\E\left[\|ru+y^*-X\|\right]&=\frac{d}{dr}\int_{x\in \R^d}\|ru+y^*-x\|d\mu \nonumber\\
        & = \int_{x\in \R^d}\frac{d}{dr}\|ru+y^*-x\|d\mu \nonumber\\
        & = \int_{x\in \R^d}\frac{\langle ru+y^*-x, u\rangle}{\|ru+y-x\|}d\mu \label{eq_exp_diff}.
    \end{align}
    Since $\mu$ is invariant under rotations centered in $y^*$, we know that a random vector $X$ drawn according to $\mu$ with $\|y^*-X\|=R$ is actually drawn according to the uniform distribution on the sphere $\SS^{d-1}_R(y^*)$, i.e., $X\sim \U(\SS_R^{d-1}(y^*))$. We evaluate \eqref{eq_exp_diff} in a fixed $r_0>0$. Let $\mu_{R}$ be the probability measure of the random variable $R=\|y^*-X\|$ and let $Z\sim \U(\SS_R^{d-1}(y^*))$. We have that
    \begin{align}
        \left.\frac{d}{dr}\E\left[\|ru+y^*-X\|\right]\right|_{r_0}& = \int_{x\in \R^d}\frac{\langle r_0u+y^*-x, u\rangle}{\|r_0u+y^*-x\|}d\mu \nonumber \\
        & = \int_0^{+\infty}\frac{\Gamma\left(\frac{d}{2}\right)}{2\pi^{\frac{d}{2}}R^{d-1}}\left(\int_{z\in\SS^{d-1}_R(y^*)}\frac{\langle r_0u+y^*-z, u\rangle}{\|r_0u+y^*-z\|}d\HH^{d-1}\right)d\mu_R,\label{eq:dobleint}
    \end{align}
    where $\HH^{d-1}$ stands for the $(d-1)$-dimensional Hausdorff measure. Now, we study the positivity of the inner integral in \eqref{eq:dobleint} by cases. In the first case, we consider $R\leq r_0$, and obtain
    \begin{equation*}
        \langle r_0u+y^*-z, u\rangle = r_0-\langle z-y^*,u\rangle \geq r_0-R\geq 0,
    \end{equation*}
    where the chain of inequalities holds tight if and only if $z = r_0u+y^*$. So we get that the inner integral in \eqref{eq:dobleint} is strictly positive when $R \in (0, r_0]$. 

    Otherwise, $R>r_0$. We define the random variable $\alpha\in[0, \pi]$ to be the angle between $z-y^*$ and $u$, and we let $\mu_\alpha$ be its associated measure. We also define the random variables $\beta(\alpha), \gamma(\alpha):[0, \pi ]\to [0, \pi ]$ to be the angles between $r_0u+y^* -z$ and $u$, and between $y^*-z$ and $r_0u+y^*-z$, respectively. It is easy to see that the three angles shape a triangle with vertices $y^*, z$, and, $r_0u+y^*$, so $\alpha+\beta(\alpha)+\gamma(\alpha)=\pi$, and therefore $\beta(\alpha)\leq \pi-\alpha$, being tight just when $\alpha=0,\pi$. Finally,

    \begin{align}
        \int_{z\in\SS^{d-1}_R(y^*)}\frac{\langle r_0u+y^*-z, u\rangle}{\|r_0u+y^*-z\|}d\HH^{d-1} & = \int_0^\pi\cos\beta(\alpha)d\mu_\alpha \label{eq:int1}\\
        &= \int_0^{\frac{\pi}{2}}(\cos\beta(\alpha)+\cos\beta(\pi-\alpha) )d\mu_\alpha\label{eq:int2}\\
        & > \int_0^{\frac{\pi}{2}}(\cos(\pi-\alpha)+\cos \alpha )d\mu_\alpha=0,\label{eq:int3}
    \end{align}
where \eqref{eq:int1} follows from scalar product formula, \eqref{eq:int2} follows from the change of variable $\alpha\mapsto \pi-\alpha$ and the symmetry of $\mu_\alpha$ with respect to $\frac{\pi}{2}$, and \eqref{eq:int3} follows from $\beta(\alpha)< \pi-\alpha$ with $\alpha\in (0,\frac{\pi}{2})$. We get that the inner integral in \eqref{eq:dobleint} is also strictly positive when $R>r_0$.
\end{proof}

\end{document}
